% Suggested insertion points:
% - end of Experimental Setup
% - start of Track B
% - appendix discussion of evaluation scope

\paragraph{Matched-scope interpretation.}
The three empirical tracks should not be read as if they were all measured on a single pooled denominator. Track~A is the primary result and uses the canonical 64-prompt suite with the first prompt excluded as warmup, yielding $N{=}63$ scored prompts across six categories. The main Track~A table remains on that denominator. When the paper needs a direct cross-track comparison, it instead uses the shared expository subset defined by canonical prompt IDs $P1$--$P5$, with $P1$ excluded as warmup so that $P2$--$P5$ are scored. We therefore treat Track~B as a scoped negative baseline and Track~C as boundary-condition evidence, and we report any direct Track~A/B/C comparison on that matched subset rather than mixing it with the canonical Track~A $N{=}63$ table.
